% Szkic artykułu pod TMLR (2026-10-27). Kompilacja: Overleaf z oficjalnym szablonem TMLR (plik tmlr.sty obok main.tex);
% bez tmlr.sty plik kompiluje się w zwykłym układzie article. Każda liczba ma źródło w README.md (tabela źródeł).
\documentclass[10pt]{article}
\IfFileExists{tmlr.sty}{\usepackage{tmlr}}{%
  \usepackage[margin=1in]{geometry}\usepackage[round]{natbib}\bibliographystyle{plainnat}}
\usepackage[utf8]{inputenc}
\usepackage[T1]{fontenc}
\usepackage{amsmath,amssymb,amsfonts}
\usepackage{booktabs}
\usepackage{multirow}
\usepackage{hyperref}
\usepackage{url}
\usepackage{xcolor}
\newcommand{\todo}[1]{\textcolor{red}{[TODO: #1]}}
\newcommand{\dep}{\operatorname{dep}}
\newcommand{\R}{\mathbb{R}}

\title{Gradient training turns transformer MLP sublayers into\\ non-normal transient amplifiers}
\author{Anonymous authors\\ Paper under double-blind review}

\begin{document}
\maketitle

\begin{abstract}
We measure the token-wise Jacobian of the MLP sublayer of vision transformers across training and ask what its structure is and what it
is for. In a 6-block detection ViT trained with 24 seeds, and in a SmallViT and a DeiT-Tiny trained on CIFAR-10, gradient training makes
this Jacobian strongly non-normal: Henrici's departure from normality rises from the random-matrix level of 0.76 to 0.88--0.94, in every
seed, without weight decay and even with random labels, where the shape is kept but the scale of the field collapses. The non-normality is
dynamically real. The Kreiss constant grows during training, and nonlinear perturbations of the network along the worst direction are
amplified by exactly the largest singular value of the linearisation (to 1\%), up to $4.5\times$ per sublayer while its spectral
radius stays at 1.1--1.6. Its structure matches the algebra of balanced amplification in cortex: the nilpotent part of the Schur form concentrates in 3--8
modes, and what the field writes is read back by it only weakly. Where the amplified mode reads and writes differs from the cortical
picture. Its input direction is aligned with the gradient of the network's readout ($30$--$50\times$ chance) and is shared by the MLPs of
successive blocks. Its output direction is read weakly and diffusely by everyone downstream, so the MLPs act as parallel amplifiers on a
shared task coordinate rather than as a feedforward chain. Ablations show that this structure is one of several equivalent solutions:
freezing the MLP at random initialisation costs at most 0.03 AP, because attention compensates. All predictions were registered with
numeric thresholds before each measurement, and we report all 79 of them, 40 of which failed.
\end{abstract}

\section{Introduction}

The MLP sublayer holds two thirds of the parameters of a transformer block. Interpretability work has read it as a key--value memory
\citep{geva2021transformer} or as a steering field acting on the residual stream \citep{mudarisov2026feedforward}, and spectral studies of
trained transformers have recently turned to the geometry of block Jacobians \citep{aubry2025transformer,fernando2026dynamics} and to
the non-normality of the attention propagator \citep{li2026transient}. What has not been measured is the linear operator the MLP
applies to a token, $J=\partial\,\mathrm{MLP}(x)/\partial x$, as it changes during training, and whether its structure has a dynamical
consequence that eigenvalues do not reveal.

Non-normal operators, whose eigenvectors are not orthogonal, are the textbook example of such hidden structure. A stable spectrum does
not bound their response in finite time: perturbations along particular directions grow transiently before they decay
\citep{trefethen1993hydrodynamic,trefethen2005spectra,schmid2007nonmodal}. The same mathematics explains subcritical transition in shear
flows, reactivity of ecological communities \citep{neubert1997alternatives}, transient amplification in directed networks
\citep{asllani2018structure}, anomalous amplification in lossy photonic media \citep{makris2014anomalous}, and, in cortex,
\emph{balanced amplification}: strong recurrent excitation stabilised by inhibition hides a feedforward chain in the Schur basis of the
connectivity matrix, which selectively and rapidly amplifies particular activity patterns \citep{murphy2009balanced,hennequin2012nonnormal,
hennequin2014optimal}, an effect now measured in primate neocortex \citep{pattadkal2023primate}.

We ask whether gradient training installs this structure in the MLPs of vision transformers \citep{dosovitskiy2021image}, and if so, what
the network uses it for. Our
contributions are:
\begin{enumerate}
\item \textbf{A replicated, pre-registered measurement.} Training makes the token-wise MLP Jacobian non-normal in 24/24 seeds of a detection
  ViT and in two CIFAR-10 architectures; the shape is installed by gradient descent regardless of the label signal, while the signal
  sets the scale (\S\ref{sec:installs}).
\item \textbf{Causal relativisation.} Freezing, parallelising and removing sublayers shows that the structure is not necessary for the task:
  attention or the MLP alone recovers the function (\S\ref{sec:compensation}).
\item \textbf{A dynamical consequence, verified nonlinearly.} Transient amplification and the Kreiss constant grow in training; the
  network itself amplifies worst-direction perturbations by the predicted singular value, while its output stays robust
  (\S\ref{sec:transient}).
\item \textbf{The structure of the amplifier.} A low-rank hidden feedforward chain with a ``valve'', originating in the weights $W_2W_1$ and
  scaled per token by the number of active GELU units (\S\ref{sec:schur}).
\item \textbf{Where it reads and writes.} The amplified mode reads the network's readout coordinate, shared across blocks, and writes
  a weak, diffuse correction; a fifth of its write in the detection model changes the token norm, which turns out to be functionally
  irrelevant for the readout (\S\ref{sec:readwrite}).
\end{enumerate}

\section{Background and related work}
\label{sec:background}

\paragraph{Non-normality.} A matrix $J$ is normal if $JJ^\top=J^\top J$. Write its Schur form $J=Q(\Lambda+N)Q^*$ with $N$ strictly upper
triangular; Henrici's departure from normality is $\|N\|_F^2=\|J\|_F^2-\sum_i|\lambda_i|^2$ \citep{henrici1962bounds}. We report the
scale-free ratio $\dep(J)=\|N\|_F/\|J\|_F\in[0,1]$; for a real Ginibre matrix of size 192 it is about $0.71$, and for the MLP Jacobian at
initialisation we measure $0.76$. For the residual update $A=I+J$ we use the one-step transient gain $G_1=\|A\|_2/\rho(A)$ with
$\rho(A)=\max_i|1+\lambda_i|$, the $k$-step gain $G_k=\|A^k\|_2/\rho(A)^k$, the numerical abscissa
$\omega=\lambda_{\max}\big((J+J^\top)/2\big)$ compared with $\max_i\mathrm{Re}\,\lambda_i$, and the Kreiss constant
$\mathcal{K}(A)=\sup_{|z|>\rho}(|z|-\rho)\,\|(zI-A)^{-1}\|_2$ evaluated on a grid \citep{trefethen2005spectra}. All four exceed their
spectral values only for non-normal $A$.

\paragraph{Balanced amplification and related effects.} In excitatory--inhibitory networks, the Schur form of the connectivity contains
large feedforward weights from ``difference'' to ``sum'' patterns; the network amplifies selected inputs quickly without slow modes
\citep{murphy2009balanced}. Hidden feedforward chains also store memory without feedback \citep{goldman2009memory,ganguli2008memory}, and
trained recurrent networks develop chain-like non-normal dynamics \citep{orhan2019improved,kerg2019nonnormal,bondanelli2020coding,
christodoulou2022regimes}. Inhibition stabilisation and the paradoxical response are related but distinct diagnostics
\citep{ozeki2009inhibitory,rubin2015stabilized,sadeh2021inhibitory}. Across fields the common object is direction-selective finite-time
gain of a spectrally stable operator \citep{mccoy2013amplification,dixon2016transient,symon2018nonnormality,baggio2020efficient}, and
controlling it improves long rollouts of learned dynamics models \citep{pervez2026controlling}; its measured size depends
on the norm, the horizon and the choice of input and readout \citep{trefethen2005spectra,baggio2020efficient}, which we therefore state
for every measurement.

\paragraph{Spectral analyses of transformers.} \citet{aubry2025transformer} relate coupling of the top singular vectors of block Jacobians
across depth to generalisation in LLMs. \citet{fernando2026dynamics} report a monotonic depth gradient from non-normal early to nearly
symmetric late blocks of the whole residual block in LLMs. \citet{li2026transient} analyse transient reserves in the attention
propagator. Activation sparsity of MLPs arises even with random labels \citep{li2023lazy}, the linearity of feed-forward blocks is learned
\citep{whipp2026linear}, and the eigenspectrum of FFN activations has been tracked in LLMs \citep{jha2026nerve}. Parallel attention/FFN
blocks have been studied in language models \citep{sonkar2023investigating}. To our knowledge no prior work measures the Jacobian of the
MLP sublayer itself across training, its Schur structure, or its alignment with the readout; our audit of the literature is summarised in
\S\ref{sec:discussion}.

\section{Setup}
\label{sec:setup}

\paragraph{Models and data.} \emph{Detection:} a SmallViT with a four-layer convolutional stem, six pre-LayerNorm transformer blocks
(width 192, MLP width 768, GELU), 600 tokens per frame and a per-token head predicting a heat logit and four box offsets, trained to
localise the tip of a surgical instrument in cataract-surgery video frames. Twenty-four seeds were trained with the same recipe
(60 epochs, 65\,640 steps); test AP@0.5 is 0.827--0.902 over seeds on 500 frames from 25 held-out videos (472 with a ground-truth box).
\emph{Classification:} DeiT-Tiny \citep{touvron2021training} (12 blocks, CLS token, patch size 4, 65 tokens) and a SmallViT-CIFAR (the same
backbone with $4\times$ bilinear upsampling and a mean-pooled linear head) trained on CIFAR-10 for 93\,780 steps, plus runs with random
labels and without weight decay. Checkpoints span initialisation ($k_0$), an early checkpoint and the end of training.
\todo{Training recipe details (optimiser, schedule, augmentation) from the training scripts; CIFAR test accuracy of all runs.}

\paragraph{The MLP field Jacobian.} For a token in state $x$ (the input of the MLP sublayer of block $b$) we compute
\begin{equation}
J=W_2\,\mathrm{diag}\big(\mathrm{GELU}'(a)\big)\,W_1\,J_{\mathrm{LN}}(x),\qquad a=W_1\,\mathrm{LN}(x)+c_1,
\end{equation}
a $192\times192$ matrix in float64, for all instrument tokens and 20 background tokens of each of 50 diagnostic frames (detection), or 16
tokens of each of 50 diagnostic images (CIFAR). Its leading singular pair $(u_1,v_1,\sigma_1)$, $Jv_1=\sigma_1u_1$, defines the direction
the field \emph{reads} ($v_1$) and the direction it \emph{writes} ($u_1$). Medians are taken over tokens and frames; blocks are reported
individually or as the middle block.

\paragraph{Pre-registration.} The study consists of six measurement experiments. Before running any job of an experiment we wrote its
predictions with numeric thresholds and a decision rule, and we wrote its evaluation script before reading any output. Thresholds were
never changed afterwards; observations made after the fact are labelled as such. Predictions are evaluated by a sign test over runs
(24/24 seeds, $p\approx10^{-7}$; 5/5 runs, $p=0.031$). Appendix~\ref{app:prereg} lists all 79 predictions and controls with their
outcome: 39 passed and 40 failed.

\section{Results}

\subsection{Training installs non-normality: shape from gradient, scale from signal}
\label{sec:installs}

\begin{table}[t]
\centering
\caption{Departure from normality $\dep(J)$ of the MLP Jacobian, medians over 24 seeds (detection).}
\label{tab:dep}
\begin{tabular}{lcccccc}
\toprule
block & 1 & 2 & 3 & 4 & 5 & 6\\
\midrule
initialisation ($k_0$) & 0.763 & 0.762 & 0.762 & 0.763 & 0.762 & 0.762\\
end of training & 0.879 & 0.909 & 0.911 & 0.916 & 0.916 & 0.886\\
\bottomrule
\end{tabular}
\end{table}

In every one of the 24 detection seeds, $\dep(J)$ of block 4 rises from 0.759--0.768 at initialisation to 0.899--0.939 at the end of
training (Table~\ref{tab:dep}). The rise happens in the first five epochs together with the task: in block 4 it goes 0.763, 0.772,
0.829, 0.902, 0.923, 0.927 at steps 0, 300, 1\,000, 3\,000, 6\,000, 10\,000 and settles at 0.916, while AP reaches 0.54 at step 1\,000
and 0.82 at step 6\,000. The depth profile is not monotonic: blocks 3--5 are the most non-normal and block 4 exceeds block 6 in every
seed, in contrast with the monotonic depth gradient reported for whole residual blocks of LLMs \citep{fernando2026dynamics}. Each token
is a saddle of the field (about 92 eigenvalues with positive and 98 with negative real part). Both CIFAR architectures replicate the
effect (middle block 0.75--0.77 to 0.92--0.93), but there the last block is the most non-normal (0.95--0.96).

Non-normality is not a signature of learning the task. Without weight decay it reaches 0.900--0.927. With random labels the detector
reaches AP 0.005--0.009 and $\dep(J)$ still rises to 0.925--0.955; on CIFAR, random labels give 0.93. What the label signal controls is the
\emph{scale} of the field: the median $\|J\|$ of block 4 goes from 3.3 to 4.7 in normal training, to 0.5 with random labels, and to 11.2
when attention is frozen; on CIFAR with random labels it falls from 3.26 to 0.41. Gradient descent gives the MLP its non-normal shape
regardless of the signal; the signal decides how large the field is.

\subsection{The structure is one of several equivalent solutions}
\label{sec:compensation}

In the trained detector the MLP builds the class separation between instrument and background tokens on the LayerNorm sphere (sum of
separation gains over blocks 0.69--1.35 for MLP sublayers versus 0.12--0.42 for attention, 24/24 seeds), while attention contracts the
token cloud by motion tangent to the sphere. Ablations retrained from scratch (three seeds each) show that this division of labour is a
property of the learned solution, not a necessity (Table~\ref{tab:ablation}). Freezing every MLP at its random initialisation keeps
$\dep(J)$ at 0.76 and costs at most 0.03 AP; attention then builds the separation (0.42--0.80). Freezing attention costs 0.02--0.08 AP;
the frozen attention becomes almost an identity (contraction product 0.92--0.97) and the MLP becomes more non-normal (0.95--0.97).
A parallel block loses nothing. Only freezing both sublayers, so that only the stem and the head learn, removes 0.25--0.43 AP. The one
component that is critical in all three systems is the MLP of the first block: removing it lowers the detector's median IoU from 0.80 to
0.24, DeiT-Tiny accuracy from 0.76 to 0.09--0.11 and SmallViT-CIFAR accuracy from 0.85 to 0.55--0.72.

\begin{table}[t]
\centering
\caption{Ablations (detection, seeds 3047 / 4047 / 5047): AP@0.5 and $\dep(J)$ of block 4 at the end of training.}
\label{tab:ablation}
\begin{tabular}{lll}
\toprule
variant & AP@0.5 & $\dep(J)$ block 4\\
\midrule
base & 0.888 / 0.885 / 0.864 & 0.93 / 0.93 / 0.90\\
parallel attention + MLP & 0.895 / 0.869 / 0.873 & 0.93 / 0.92 / 0.92\\
MLP frozen at initialisation & 0.866 / 0.857 / 0.884 & 0.76\\
attention frozen at initialisation & 0.861 / 0.861 / 0.786 & 0.95 / 0.97 / 0.95\\
both frozen (stem and head learn) & 0.463 / 0.491 / 0.618 & 0.76\\
no weight decay & 0.886 / 0.873 / 0.865 & 0.92 / 0.93 / 0.90\\
random labels & 0.005 / 0.005 / 0.009 & 0.92 / 0.96 / 0.94\\
\bottomrule
\end{tabular}
\end{table}

\subsection{Non-normality has a dynamical consequence}
\label{sec:transient}

We measured five runs (three detection seeds, SmallViT-CIFAR, DeiT-Tiny) at initialisation, early and at the end of training. In all five,
the Kreiss constant of the middle block grows from 2.0--2.5 to 3.0--5.1, and the gap between the numerical abscissa and the spectral
abscissa grows 2--10$\times$. The effect is strong in detection, where the one-step gain $G_1$ reaches 2.15 / 2.59 / 1.66 and the
largest $k$-step gain ($k\le6$) 5.1 / 5.3 / 3.7, and weak in classification ($G_1$ 1.30 for SmallViT-CIFAR and 1.14 for DeiT-Tiny, against
1.07--1.13 at initialisation).

The linearisation is faithful. We perturbed a token at the input of the MLP sublayer by $\delta=0.01$ of its norm along the leading right
singular vector and propagated it through the nonlinear network. The measured amplification equals $\sigma_{\max}$ to within 1\% in all
runs (measured / predicted: 3.21 / 3.22, 4.47 / 4.51, 2.06 / 2.07, 1.47 / 1.48, 1.21 / 1.21). The worst direction changes the output
2.6--7$\times$ more than a random direction of the same size. The output itself is robust. A worst-direction perturbation of $\delta=0.03$
changes the IoU of the detected box by 0.005--0.013 (CIFAR margin: 0.07--0.11), no block dominates the sensitivity, and composing the
blocks does not compound the gain in detection: the product of six residual updates amplifies less (0.97--1.21) than the median single
block (1.65--1.76), because the worst directions of consecutive blocks are not aligned. Background tokens are amplified more than
instrument tokens (for one seed $G_1$ 2.76 versus 1.44).

\subsection{A hidden low-rank feedforward chain with a valve}
\label{sec:schur}

\begin{table}[t]
\centering
\caption{Schur structure of the MLP Jacobian, middle block, initialisation $\to$ end of training (five runs).}
\label{tab:schur}
\begin{tabular}{lll}
\toprule
quantity & initialisation & end of training\\
\midrule
share of $\|N\|_F^2$ in the top 5 singular values of $N$ ($c_5$) & 0.15--0.16 & 0.63--0.74\\
effective rank of $N$ ($\mathrm{PR}_N$) & 60--63 & 3.3--8.3\\
valve $\|Ju_1\|/\sigma_1$ & 0.39--0.45 & 0.18--0.36\\
$\dep(W_2W_1)$ (weights only) & 0.73--0.74 & 0.85--0.90\\
$|\cos(w_{1,i},w_{2,i})|$, gate-weighted mean over units & 0.053--0.059 & 0.10--0.16\\
\bottomrule
\end{tabular}
\end{table}

The Schur form makes the structure explicit (Table~\ref{tab:schur}). We computed it without SciPy by eigendecomposition, QR of the
eigenvectors and $T=R\Lambda R^{-1}$, and took $N$ as the strictly upper part of $T$. At initialisation the nilpotent part is spread over
60--63 effective dimensions, as for a random matrix. After training, 63--74\% of $\|N\|_F^2$ sits in five modes and its effective rank is
3.3--8.3, in all five runs: a low-rank hidden feedforward chain in the sense of \citet{murphy2009balanced} and \citet{goldman2009memory}.
The leading mode has a valve: what the field writes, $u_1$, is mapped by the field with only 0.18--0.36 of the gain $\sigma_1$, against
0.39--0.45 at initialisation.

The non-normality originates in the weights. The product $W_2W_1$ alone reaches $\dep=0.85$--$0.90$ (0.73 at initialisation), and
inserting the mean gate, $W_2\,\mathrm{diag}(\bar g)\,W_1$, matches $\dep(J)$ to within 0.02 in detection (0.05--0.06 on CIFAR). The gate
controls the size of the amplification per token: across tokens of the middle block, the number of active units ($0.1<\mathrm{GELU}'<0.9$)
correlates with $G_1$ at 0.41--0.84 at the end of training ($-0.38$ to $+0.11$ at initialisation). On instrument tokens the MLP is almost
switched off: they have fewer active units than background tokens (73--103 versus 127--162) and a mean gate derivative near zero, which
explains their smaller transient gain. Individual units become partly reciprocal (they write closer to the direction they read,
$|\cos(w_{1,i},w_{2,i})|$ 0.10--0.16 against 0.057 for random weights) while the non-normality of the whole field grows; a minority of
reciprocal units coexists with a majority forming the chain.

\subsection{What the amplified mode reads, and who reads what it writes}
\label{sec:readwrite}

\paragraph{It reads the readout coordinate, shared across blocks.} Let $r$ be the gradient of the token's own output (detection: its heat
logit; CIFAR: the class margin) with respect to the token state after block $b$. In detection the \emph{input} direction $v_1$ of the
amplified mode is aligned with $r$: $\cos^2(v_1,r)=0.06$--$0.11$ in the middle block, 30--50$\times$ the random level of 0.002. Its
\emph{output} direction is not ($\cos^2(u_1,r)=0.003$--$0.008$). This reverses our pre-registered prediction, which, following the cortical
picture, expected the written direction to be aligned with the readout. In the subspace spanned by the gradients of all five outputs of
the token, $v_1$ carries 0.14--0.17 of its energy (random directions 0.022--0.023). The input directions of successive blocks coincide:
$\cos^2(v_1^{b},v_1^{b+1})=0.09$--$0.21$, against 0.003 at initialisation and 0.0052 for random directions.

\paragraph{The write is not a message to a specific reader.} We tested the candidate readers of $u_1$ directly. The detection head does
not read it: its energy in the box-output subspace is 0.021--0.039 (chance $4/192=0.021$) and in the five-output subspace 0.030--0.046
(random directions 0.022--0.023); only in the last block does it enter the head subspace (4--6$\times$ chance). The next MLP does not read it with its leading direction
($\cos^2(u_1^b,v_1^{b+1})=0.008$--$0.021$). For attention in block $b+1$ we computed the linear read maps $W_{q,k,v}J_{\mathrm{LN}}$ and the
gain of $u_1$ relative to random directions of the same token. Attention reads $u_1$ through queries/keys $2.3$--$4.0\times$ chance (DeiT
3.6) and through values $2.1$--$2.4\times$ (DeiT 4.5), with a ceiling of 16--26 for the best direction of each reader; the next MLP reads it
with 1.7--3.4$\times$ and the head with 1.8--2.4$\times$ in detection (4.3--4.6$\times$ on CIFAR). Nonlinearly, writing along $u_1$ in one token changes the attention output of the
\emph{other} tokens 1.5--1.8$\times$ more than a random write (DeiT 2.3), below our pre-registered 2$\times$ threshold, and it shifts the
balance between routing and content exactly as a random write does. The write is not lost: the sensitivity of the network output along
$u_1$ is 1.3--2.2$\times$ chance (5/5 runs); in detection it is 2--3$\times$ smaller than along the read directions, on CIFAR the two are
comparable. By contrast, the read direction of the next block is read strongly by everyone (next MLP 9--12$\times$ and head
21--22$\times$ in detection).

\paragraph{On CIFAR the write goes to the readout.} With a global readout the picture changes: $u_1$ carries 0.17--0.20 of its energy in the
subspace of the ten logit gradients in the middle block (random directions 0.048--0.050), 0.82 in the last block of SmallViT-CIFAR, and the
leading direction of the next MLP picks it up above our threshold ($\cos^2(u_1^b,v_1^{b+1})=0.017$--$0.021$, 3--4$\times$ chance). The
direction of alignment with the task therefore depends on the type of readout: per-token in detection, global in classification.

\paragraph{A fifth of the write changes the token norm, without functional effect.} In the detection model 20--25\% of the energy of $u_1$
lies in the null space of the next LayerNorm (8\% for SmallViT-CIFAR; 0.7\% for DeiT-Tiny, the same as for random directions). Decomposing this null space,
virtually all of it is the radial direction of the centred token (energy 0.197 / 0.253 / 0.240 in the three detection seeds), not the
mean direction ($\le0.001$); it is learned (0.001--0.005 at initialisation). When a token is pushed towards higher readout along $v_1$, the
MLP shrinks its norm in 66--71\% of tokens (oriented radial cosine $-0.23$ to $-0.25$), although its actual update inflates tokens by
16--21\% per block. Because every downstream reader sits behind LayerNorm, the norm only matters through the relative size of later updates.
We measured this counterfactually by removing the radial component of the MLP response and propagating through the rest of the network:
it accounts for at most 0.14\% of the change in the token's readout ($\le1.5\%$ at $\delta=0.1$), against 18--24\% carried by the
perpendicular part of the response and 77--80\% by the identity path. Our pre-registered hypothesis that the norm acts as a divisive,
stabilising channel \citep{carandini2012normalization} is rejected; the norm channel exists but is irrelevant for the readout.

\section{Discussion}
\label{sec:discussion}

\paragraph{What the MLP is doing.} Taken together, the measurements describe each MLP sublayer as a non-normal amplifier connected to a
shared bus, in the sense in which the residual stream is a shared communication channel \citep{elhage2021mathematical}. All blocks read the same task coordinate, the one the head reads; each amplifies selected perturbations of it transiently,
by up to $4.5\times$ per sublayer; and each writes the result into directions that the field itself and the following readers pick up only
weakly and diffusely. This is closer to a set of parallel differential amplifiers than to the feedforward chain of balanced amplification,
in which the amplified pattern is passed on and read downstream \citep{murphy2009balanced,goldman2009memory}. It also differs from the
coupling of top singular vectors across blocks reported for LLMs \citep{aubry2025transformer}: we find coupling of what blocks read, not of
what one block writes and the next reads. In the key--value reading of MLPs \citep{geva2021transformer}, our result says that the keys of
the most amplified mode are tuned to the readout coordinate while its values are spread; the amplification is an internal property of the
representation rather than a direct contribution to the output, which stays robust.

\paragraph{Shape and scale.} The non-normal shape appears with random labels, without weight decay and at the same early stage of training
in every seed, while its scale follows the label signal. This places it with the cases where non-normality is created by optimisation
rather than given by design, as in trained recurrent networks \citep{orhan2019improved} and in the optimal control of cortical dynamics
\citep{hennequin2014optimal}. Because freezing the MLP at initialisation hardly hurts the detector, the structure should be read as the
generic solution gradient descent finds for this sublayer, not as a mechanism the task requires.

\paragraph{The physics analogy, stated precisely.} The shared mathematics is that of a spectrally stable operator whose finite-time gain is
set by its singular vectors, pseudospectra and Kreiss constant \citep{trefethen2005spectra,schmid2007nonmodal,symon2018nonnormality}.
The MLP Jacobian has the algebra of balanced amplification (a low-rank Schur chain, a valve, stable eigenvalues) but not its
excitatory--inhibitory origin, and its alignment with the task is on the input side. In the language of non-normal photonic amplifiers,
the optimal excitation of the medium lies along the readout direction \citep{makris2014anomalous}, and the valve is the analogue of
directional amplification \citep{midya2022topological,ramos2021topological}. The norm channel of \S\ref{sec:readwrite} is reminiscent of
the high-norm tokens triggered by the leading right singular vector of a module's linearisation in large models
\citep{wang2025demystifying}, although in our detector it has no measurable effect on the readout. Measures of
non-normal amplification depend on norm, horizon, input and readout \citep{baggio2020efficient}; we report all four, and we validate every
linear prediction on the nonlinear network.

\paragraph{Novelty.} A literature audit (two systematic searches, 21 and 134 papers) found no measurement of the MLP sublayer Jacobian in
trained transformers with Schur structure, Kreiss constant or readout alignment; the closest works analyse attention \citep{li2026transient}
or whole residual blocks \citep{fernando2026dynamics,aubry2025transformer} in language models.

\section{Limitations}

The models are small (width 192, 6 and 12 blocks) and the tasks are two; whether the structure persists in large pretrained ViTs is the
most important open question, and methods that infer non-normal amplification from activations alone \citep{saiprasad2026inferring}
could extend the measurement to models whose Jacobians are too costly to compute per token. The mechanistic measurements (\S\ref{sec:transient}--\ref{sec:readwrite}) use five runs, three of them
seeds of one detector and one seed per CIFAR architecture, which supports the direction of effects but not tight intervals. The analysis
is local, around the trained trajectory, although every linear prediction we rely on was checked nonlinearly. Readouts are single outputs
per token (detection heat logit) or the class margin; other readouts could change the alignment results. We have not found an
intervention that removes the non-normal structure while keeping the learned function, so our causal statements are limited to the
ablations of \S\ref{sec:compensation}. \todo{Figures: departure across training; Schur concentration; nonlinear gain vs.\ $\sigma_{\max}$;
read/write alignment diagram.}

\section{Conclusion}

Gradient training gives the MLP sublayer of vision transformers a non-normal, low-rank, valved Jacobian whose amplified mode reads the
network's readout coordinate. The amplification is real and predicted to 1\% by the singular values, yet the output stays robust and the
written signal is received only diffusely. The structure is the generic product of gradient descent on this sublayer, its scale is set by
the task signal, and the network can do without it.

\bibliography{references}
\IfFileExists{tmlr.sty}{\bibliographystyle{tmlr}}{}

\appendix
\section{Pre-registered predictions and their outcomes}
\label{app:prereg}
\todo{Full table of all predictions per experiment with thresholds and outcomes, generated from the OCENA\_PREDYKCJI files of the six
experiments (see README).}

\end{document}
